\begin{pastexam}[label={px:ncku-110-ee-ds-06}]{成大電機所 110 學年度 第 6 題}
（問答）

Describe three key points for designing a good hashing function. \hfill (15 \%)
\end{pastexam}

\begin{pastsoln}
好的雜湊函數（資結書第~12.1~節）要注意：

\begin{enumerate}
\item \textbf{算得快}。 每次搜尋、插入、刪除都要算一次雜湊值，雜湊函數本身應該是 $O(1)$（或和鍵的長度成
正比），只用幾個算術或位元運算。 算得慢就失去雜湊表「平均 $O(1)$」的意義。
\item \textbf{分布均勻、碰撞少}。 鍵應該盡量平均地分到 $m$ 個位置，每個位置的機率接近 $1/m$（簡單均勻
雜湊），這樣串列或探測序列的平均長度才會是 $O(1 + \alpha)$。 不要讓實際資料常見的規律（例如都是偶數、
都是某個數的倍數、字串前綴相同）集中到少數幾個位置。
\item \textbf{用到鍵的全部資訊}。 雜湊值應該和鍵的每一部分都有關，例如字串的每個字元都要參與計算
（$h(s) = (\sum s_i\, b^{\,i}) \bmod m$），不要只取前幾個字元或只取數字的低位元；否則只在沒用到的部分
不同的鍵全部會撞在一起。
\end{enumerate}

常見的具體做法也和這三點有關：除法 $h(k) = k \bmod m$ 的 $m$ 取不接近 $2$ 的次方的質數（$m = 2^p$
的話只用到 $k$ 的最低 $p$ 個位元，違反第 3 點；$m$ 和鍵常見的間隔有公因數時，鍵只會落在少數位置，違反
第 2 點）。 另外，雜湊函數必須是確定的：同一個鍵每次都要算出同一個值，否則存進去就找不回來。
\end{pastsoln}
